\chapter{Práctica: Grafos en Haskell}

\section{Objetivos}
Al finalizar esta práctica, el estudiante será capaz de:
\begin{itemize}
  \item Implementar representaciones de grafos en Haskell.
  \item Implementar algoritmos de recorrido y búsqueda.
  \item Implementar algoritmos de caminos mínimos.
\end{itemize}

\section{Ejercicios de Implementación}

\subsection{Representación y Recorridos}
\begin{enumerate}
  \item Implementar una función que convierta una matriz de adyacencia a una lista de adyacencia.
  \item Implementar DFS y BFS usando colas y pilas explícitas.
  \item Implementar una función que verifique si un grafo es conexo.
  \item Implementar una función que encuentre todas las componentes conexas.
\end{enumerate}

\subsection{Algoritmos Avanzados}
\begin{enumerate}
  \item Implementar el algoritmo de Dijkstra usando una cola de prioridad.
  \item Implementar el algoritmo de Prim.
  \item Implementar el algoritmo de Kruskal usando Union-Find.
  \item Implementar una función que encuentre un ciclo euleriano.
\end{enumerate}

\section{Proyecto}
Desarrollar una librería en Haskell para teoría de grafos con las siguientes funcionalidades:
\begin{itemize}
  \item Representación de grafos (no dirigidos, dirigidos, ponderados).
  \item Algoritmos de recorrido (DFS, BFS).
  \item Algoritmos de caminos mínimos (Dijkstra, Floyd-Warshall).
  \item Algoritmos de árboles generadores (Prim, Kruskal).
  \item Verificación de propiedades (conexidad, bipartismo, planaridad).
\end{itemize}